\title{Large Language Models Can Automatically Engineer Features for Few-Shot Tabular Learning}

\begin{abstract}
Large Language Models (LLMs) have exhibited exceptional proficiency in addressing complex and novel reasoning challenges, positioning them as promising tools for tabular learning—a domain integral to numerous practical applications. This work introduces \textsf{FeatLLM}, a new in-context learning approach that leverages LLMs as automated feature engineers, transforming raw input data into feature sets specifically crafted to enhance tabular prediction tasks. The resulting features are subsequently employed to estimate class probabilities using simple downstream models, such as linear regression, enabling robust performance even in few-shot learning scenarios. Notably, \textsf{FeatLLM} relies solely on this straightforward predictive model utilizing the discovered features during the inference phase. In contrast to prior LLM-driven methods, \textsf{FeatLLM} obviates the necessity of querying the LLM for every input during inference. Additionally, it requires only API-level interaction with the LLM and effectively navigates limitations imposed by prompt sizes. Experiments spanning a diverse array of tabular datasets demonstrate that \textsf{FeatLLM} produces superior rule sets and achieves significantly better accuracy—by an average of 10\%—over state-of-the-art methods like TabLLM and STUNT.
\end{abstract}

\section{Introduction}
Large Language Models (LLMs) have, in recent years, demonstrated significant aptitude in responding to novel tasks without the need for domain-specific fine-tuning~\cite{creswell2022selection,imani2023mathprompter,taylor2022galactica}. Their training on vast and varied textual data grants LLMs extensive background knowledge, often substituting for specialized domain expertise in several contexts~\cite{Genomes2021,singhal2023large,wilcox2003role}, thereby supporting automation across applications such as dialogue evaluation~\cite{finch2023leveraging} and automated code correction~\cite{fan2023automated}. By including a combination of task descriptions and a handful of examples within prompts, LLMs can grasp the underlying context and focus on salient attributes and samples~\cite{brown2020language,zhao2023survey}. This ability underlines the prospect of LLMs functioning as proficient feature engineers by analyzing and extracting relevant aspects of the data.

Applying the intrinsic reasoning and background knowledge of LLMs to tabular learning has yielded promising results. For example, knowledge about heightened disease risk associated with advanced age can directly inform disease prediction models. Building on this insight, contemporary research represents tasks and feature sets as natural language statements, serializes tabular data, and delivers this information to LLMs~\cite{dinh2022lift,wang2023anypredict}. Subsequent approaches include in-context learning~\cite{nam2023semi} or efficient fine-tuning methodologies~\cite{dinh2022lift,hegselmann2023tabllm}. Leveraging LLMs' embedded prior knowledge has shown substantial benefits, especially surpassing classic tabular learning strategies in settings where labeled data is limited~\cite{hegselmann2023tabllm}.

Existing methods for applying LLMs to tabular learning are hindered by several notable issues. For direct prediction tasks, they necessitate running at least one inference per data point, leading to substantial computational overhead. Achieving high levels of accuracy usually demands fine-tuning the LLM, which is impractical for models without accessible training frameworks or publicly available tuning interfaces—an increasing problem as many state-of-the-art LLMs now restrict access to inference-only APIs~\cite{anil2023palm,openai2023gpt4}. Moreover, these techniques often perform poorly or become unusable as prompt length increases~\cite{liu2023lost}; in scenarios where tabular datasets have many features, prompt serialization leads to excessively long text inputs, which either exceed the input limits or result in reduced performance. Collectively, these limitations create significant barriers for deploying LLM-driven tabular models in practice.

In response, our methodology reconceptualizes the application of LLMs, emphasizing their utility in feature engineering rather than solely end-to-end prediction. We present a new in-context learning paradigm, \textsf{FeatLLM}, which focuses on discerning the “criteria” used by LLMs for making predictions. For example, in disease classification, the LLM can articulate explicit rules or conditions based on features that are indicative of the disease. These identified criteria are then translated into engineered features, replacing the original variables in downstream tasks. This redirection confines the LLM’s involvement to feature creation, simplifying the subsequent modeling process—since, after discovering features, inference can be performed without sophisticated machine learning algorithms, thus reducing latency. By utilizing in-context learning, we can generate high-quality features from the LLM using prompt-based demonstrations, completely bypassing model retraining. Additionally, by incorporating ensemble techniques such as feature bagging~\cite{breiman2001random}, we can alleviate problems arising from oversized prompts.

\textsf{FeatLLM} decomposes the process of feature engineering from input prompts into two sequential tasks: (1) interpreting the problem to deduce connections between features and targets; and (2) utilizing insights gained from this analysis, along with few-shot examples, to formulate discriminative rules that can separate classes. This systematic reasoning allows LLMs to disregard irrelevant features when extracting rules. The resulting rules are executed over the dataset (interpreted as program code), generating binary features that indicate, for each sample, adherence to individual rules. Subsequently, a simple model is trained to predict class probabilities using these derived binary features. The framework’s resilience is further enhanced through an ensemble technique, in which feature generation is repeated alongside feature bagging. By carefully managing the number of features and sampled data in each bagging iteration, \textsf{FeatLLM} effectively controls prompt size constraints. Because \textsf{FeatLLM} requires no additional training and entails minimal inference overhead, it is well-suited for practical deployment across a range of applications.

We empirically assess \textsf{FeatLLM} on 13 distinct tabular datasets under limited data conditions, demonstrating its consistently strong and robust outcomes. The results indicate that LLMs are adept at leveraging both prior and dataset-specific knowledge to construct informative features. Our method consistently surpasses leading few-shot learning benchmarks in diverse scenarios. An anonymized implementation is provided at~\texttt{https://github.com/Sungwon-Han/FeatLLM}.

\section{Related Work}

\subsection{Few-Shot Learning with Tabular Data} The development of few-shot learning methodologies has made it possible to obtain robust, generalizable representations even when labeled data is scarce~\cite{chen2018closer}. While early work concentrated primarily on the image domain, recent studies have established that few-shot techniques are versatile, showing effectiveness with text, audio, and, increasingly, tabular data as well~\cite{majumder2022few,nam2023stunt,schick2021exploiting}. Nevertheless, learning from a limited set of labeled instances increases susceptibility to capturing incidental or spurious correlations~\cite{bartlett2021deep}. In response, contemporary strategies often supplement available labeled examples with external sources of information or introduce relevant domain-specific prior knowledge to instill suitable inductive biases into the learning process. For instance, researchers may utilize substantial quantities of accessible unlabeled tabular datasets, applying strategic augmentations within a semi-supervised learning paradigm~\cite{ucar2021subtab,yoon2020vime}, or pursue favorable model initializations via unsupervised pre-training unrelated to any particular downstream task~\cite{somepalli2021saint}. Unsupervised learning objectives in this context commonly involve data corruption or masking followed by a reconstruction loss~\cite{arik2021tabnet,majmundar2022met}, or leverage data augmentation to implement contrastive learning objectives~\cite{bahri2022scarf,chen2020simple}. However, due to the intrinsically heterogeneous characteristics of tabular datasets, universal augmentation schemes remain elusive, and such approaches have shown limited effectiveness for few-shot learning scenarios~\cite{nam2023stunt}.

More recent advancements advocate pre-training models across an array of diverse tabular datasets, even those unconnected to the target application, and subsequently transferring the acquired knowledge~\cite{zhang2023generative,levin2022transfer,zhu2023xtab}. One illustrative example is TransTab~\cite{wang2022transtab}, where transformer models are trained to exploit not just table values but also semantic cues from column names, fostering a deeper understanding of inter-column relationships. The abstractions learned from varied columns and tables have enabled the deployment of such models in zero-shot or few-shot contexts for previously unseen tabular tasks. In another direction, TabPFN~\cite{hollmann2022tabpfn} employs synthetic data generation using mixed distributions for pre-training a Bayesian neural network; once pre-trained, the model predicts on new tasks directly, relying only on the supplied training and test data of the target, obviating further retraining. Prior exposure to heterogeneous datasets allows these models to surpass conventional tree-based learners, yielding demonstrable benefits in few-shot scenarios.

\subsection{Language-Interfaced Tabular Learning} A complementary strategy for integrating prior knowledge involves employing large language models (LLMs)~\cite{anil2023palm,openai2023gpt4}. LLMs, which have been pretrained on massive and varied text collections, exhibit strong generalization abilities, often performing well on novel and unseen tasks across different domains~\cite{brown2020language}. Recent research has sought to leverage the substantial repository of prior information embedded in LLMs to advance tabular data modeling. A common approach entails transforming tabular inputs into textual format and embedding this structured information, alongside feature and task specifications, within the LLM's input prompt~\cite{dinh2022lift,hegselmann2023tabllm,wang2023anypredict}. Through such prompt engineering, LLMs can infer interdependencies between tasks and input features for downstream inference. Approaches within this line of work include both parameter-efficient fine-tuning methods such as LoRA~\cite{hu2021lora} and IA3~\cite{liu2022few}, as well as in-context learning strategies that involve supplying the prompt with illustrative few-shot examples in lieu of explicit model training~\cite{dinh2022lift,hegselmann2023tabllm,nam2023semi,slack2023tablet}.

Although prior methods have demonstrated efficacy in enhancing few-shot tabular learning, their reliance on routing inference through the LLM at every prediction step introduces complications for deployment in industrial contexts. This work seeks to move beyond the traditional paradigm where each data point’s output is determined in an end-to-end manner by the LLM as a black box. The approach instead utilizes the LLM to deduce explicit conditions or rules representative of each output class. \textsf{FeatLLM} subsequently converts these rules into executable program code, generating new features that facilitate inference independent of the LLM. By harnessing in-context learning, the system draws upon pre-existing knowledge as well as insights derived from the few-shot instances to produce meaningful rules.

\section{Methods}
\textsf{FeatLLM} is architected to extract the underlying “rules”—the logical criteria that specify membership in each answer class—from a limited set of example inputs, as illustrated in Figure~\ref{fig:main_model}. These extracted rules are formulated with assistance from an LLM and are then used to create a suite of binary features, each indicating whether a particular rule holds for a given data point. This collection of binary features is subjected to a dot product with a vector of non-negative, learnable weights, which serves to compute the predicted probability for each class. Through training, the model identifies the influence each binary feature has on the classification outcome (see Section~\ref{Sec:training}). To enhance robustness, this methodology is performed iteratively with bagging, and the resultant predictions are aggregated using an ensemble approach. The subsequent sections detail the problem formulation as well as each component of the pipeline.

\vspace*{-0.12in}
\paragraph{Problem formulation.} Consider a tabular dataset consisting of $N$ annotated examples, represented as $\mathcal{D} = \{(\mathbf{x}^i, \mathbf{y}^i)\}_{i=1}^{N}$. Each sample $\mathbf{x}^i$ is a $d$-dimensional vector capturing $d$ distinct input features, and each feature is accompanied by a descriptive name in natural language, indicated as $F = \{f_j\}_{j=1}^{d}$. Their detailed definitions are made available separately. In the context of classification, the target $\mathbf{y}^i$ is drawn from a fixed set of possible classes. For $k$-shot learning scenarios, we select a random subset of $k$ ($k<N$) labeled examples from $\mathcal{D}$ to use for model training.

\subsection{Prompt Design for Extracting Rules}
\label{Sec:prompt_design}
To facilitate the extraction of rules via LLMs in a manner that aligns closely with rigorous reasoning processes, we design prompts that emulate the approach a domain expert might take when engaging with tabular data analysis. Our prompt structure encompasses three core elements (see Fig.~\ref{fig:prompt_for_rule} and Appendix~\ref{sec:example_rule}):

\vspace*{-0.12in}
\paragraph{Basic information description.} This section presents the foundational details necessary for problem resolution (refer to the orange-marked content in Fig.~\ref{fig:prompt_for_rule}). Specifically, it comprises: a statement outlining the task along with information on labels; natural language descriptions and type specifications of each feature; and illustrative training instances. The task prompt is phrased as a direct question (for instance, ``Does this patient's myocardial infarction complications data indicate chronic heart failure? Yes or no?''; additional examples are available in Appendix~\ref{sec:task_desc}). Feature descriptions specify the datatype, and may further provide a concise definition; for categorical attributes, representative category values may also be included. For the purpose of demonstration, select training instances are formatted as text strings along with their true labels. This text serialization adopts the style employed by prior research~\cite{dinh2022lift,hegselmann2023tabllm}:

\begin{align}
&\text{Serialize}(\mathbf{x}^i,\mathbf{y}^i, F) = \nonumber \\ 
&``\ f_1\ \text{is}\ \mathbf{x}^i_1.\ ... \ f_d\  \text{is}\  \mathbf{x}^i_d.\ \ \text{Answer:}\  \mathbf{y}^i\ ”
\end{align}

\vspace*{-0.12in}
\paragraph{Reasoning instruction.} Our objective is to improve the reasoning abilities of the LLM by offering explicit reasoning instructions (highlighted in blue in Fig.~\ref{fig:prompt_for_rule}). The reasoning instruction commences with an opening line analogous to the chain-of-thought methodology~\cite{wei2022chain}. Following this, we decompose the inference mechanism of the LLM into two sequential phases. During the first phase, the model is tasked with deducing the potential causal dynamics or tendencies that exist between features and the task description. This step occurs independently of any provided demonstrations, leveraging only broad domain knowledge or intuitive reasoning, thus encouraging the LLM to utilize its background knowledge pertaining to the given problem. In the second phase, the LLM incorporates insights from both illustrative demonstrations and its prior causal analysis to infer class-specific rules. Adopting this bifurcated approach mitigates the risk of the model latching onto misleading correlations in irrelevant columns and supports prioritizing features of greater relevance. To ensure a balance and prevent the extraction of either too few (which may underfit) or too many rules (which could unnecessarily extend the prompt), we prescribe a set number of rules to be generated per class (specifically, 10)\footnote{For further examination of rule quantity, see Section~\ref{sec:analysis}.}. Moreover, when constructing these rules, the LLM is directed to reference the feature types from the basic information segment to avert mismatches in feature-rule correspondence.

\vspace*{-0.12in}
\paragraph{Response instruction.} To streamline the post-inference processing and application of generated rules, we steer the LLM’s response organization via explicit instructions (marked in yellow in Fig.~\ref{fig:prompt_for_rule}). The prompt stipulates detailed formatting expectations for both the overall response of each answer class and for each individual rule. By doing so, the LLM is able to tailor the output format according to the nature of the corresponding feature type. Absent additional directives, the LLM is permitted to combine rules with logical connectors such as AND or OR. Illustrative results can be found in Appendix~\ref{sec:outcome-rule}.

\subsection{Inferring Class Likelihood via Rules}
\label{Sec:training}

\paragraph{Parsing rules for feature generation.} In this step, we transform the rules obtained from the previous phase into additional binary features. Each of these features corresponds to a particular class and indicates whether a given sample satisfies the set of rules specific to that class (taking values '0' or '1'). Such features assist in assessing the probability that a sample belongs to a specific class. Nonetheless, as the rules from the LLM are produced in natural language, an automated parsing approach is necessary for effective feature extraction. Despite imposing formatting guidelines during rule construction to streamline parsing, outputs from the LLM may still exhibit variations in syntax and introduce irregularities~\cite{kovavc2023large}. For example, the LLM might replace operators like ``$>$” or ``$<$” with their verbal equivalents such as ``greater than'' or ``less than,'' complicating the parsing process.

To overcome the difficulties posed by noisy or variable textual rules, rather than developing intricate parsing logic, we utilize the LLM itself as the parser. The LLM excels at comprehending underlying meaning even amid syntactic discrepancies and is adept at synthesizing executable code from instructions, owing to its exposure to code in its training data. To exploit these capabilities, our approach involves constructing prompts that include the function signature, explanations of inputs and outputs, and the parsed rules, and then providing these prompts to the LLM. Illustrative examples of these prompts and their generated outputs can be found in Figures~\ref{fig:prompt_for_parsing} and~\ref{sec:example_parsing}-\ref{sec:example_parse_outcome} in the Appendix. The resulting code snippets are then run using Python's \texttt{exec()} function to facilitate feature conversion.

\vspace*{-0.12in}
\paragraph{Inferring class likelihood.} The newly introduced binary features for each class enable us to estimate how likely a sample is to belong to each class, initially by tallying the number of class-specific rules a sample meets (i.e., summing the relevant binary features per class). However, the discriminative power of each rule may vary, making it necessary to determine their individual contributions using the training data. For this purpose, we adopt a standard machine learning (ML) model to quantify the significance of each feature within a class. A straightforward approach is to use a linear model without a bias term. Formally, for a sample $\mathbf{x}^i$ and class $k$, the binary feature vector is represented as $\mathbf{z}_k^i \in \{0, 1\}^R$, where $R$ is the number of class-specific rules and $c$ indicates the number of classes. The class probability $\mathbf{p}^i$ is calculated as follows:

\begin{align}
\text{logit}_k^i &= \max(\mathbf{w}_k, 0) \cdot \mathbf{z}_k^i, \nonumber \\
\mathbf{p}^i &= \text{Softmax}({[\text{logit}_{1}^i, ..., \text{logit}_{c}^i]}). \label{eq:infer_prob}
\end{align}

In this setup, $\mathbf{w}_k$ denotes the learnable parameters within the linear model, reflecting the relative contribution of each feature. By constraining these weights to a positive domain, we guarantee that features derived from the LLM will not diminish the predicted probability of a target class during model training. The resulting logits are transformed into probability estimates over classes via the softmax operation. To determine the optimal values for $\mathbf{w}_k$, we minimize the cross-entropy loss utilizing a limited set of training examples in a few-shot manner.

\vspace*{-0.12in}
\paragraph{Ensembling with bagging.} The procedure is iterated multiple times, yielding several models that collectively inform the final decision through ensembling. Specifically, the ensemble prediction aggregates the individual class probabilities $\mathbf{p}^i$ computed using the distinct weight vectors $\{\mathbf{w}[t]\}_{t=1}^T$ obtained across $T$ independent runs, as formalized in Eq.~\ref{eq:infer_prob}.

To foster diversity in the ensemble’s rule generation, the temperature parameter during LLM inference is set to a sufficiently high value. Furthermore, the few-shot examples are presented in random sequences, recognizing that the arrangement of these demonstrations can influence the LLM’s responses~\cite{lu2022fantastically}\footnote{Refer to Appendix~\ref{sec:diversity} for an analysis of rule diversity.}. In addition, bagging is employed by subsampling features or training instances for each trial, which is especially advantageous when handling datasets with large sample sizes or high feature dimensionality. This ensemble technique brings two primary benefits. First, should the LLM introduce erroneous rules in certain iterations, those errors are mitigated by accurate rules discovered in other runs, thereby improving framework reliability. This benefit aligns with the principle of self-consistency in LLMs~\cite{wang2022self}, whereby rules consistently produced in multiple attempts have increased credibility. Second, the ensemble strategy counters the constraint of LLM prompt length: when tabular datasets exceed prompt size limits, bagging effectively downsizes the problem, helping to preserve the model’s effectiveness. Consequently, though any single model may miss specific feature or instance relationships, the aggregate of multiple weak models constructed over different subsamples helps counterbalance partial coverage and enhances generalization.

\section{Experiments}
We assess the performance of \textsf{FeatLLM} across a variety of tabular datasets and perform several analytical studies to understand the underlying mechanisms of our proposed framework. Due to page limitations, supplementary experiments—including those involving alternative LLMs, qualitative evaluations, and further investigations—are documented in the Appendix.

\subsection{Performance Evaluation}
\paragraph{Datasets.}
Our study makes use of 13 datasets that cover binary and multi-class classification scenarios: (1) Adult~\cite{asuncion2007uci}, focused on determining if an individual’s yearly income exceeds \$50,000; (2) Bank~\cite{moro2014data}, predicting a client’s likelihood of subscribing to a term deposit; (3) Blood~\cite{yeh2009knowledge}, classifying whether a blood donor is likely to return; (4) Car~\cite{kadra2021well}, assessing the overall quality of cars; (5) Communities~\cite{misc_communities_and_crime_183}, estimating the crime rate in a specified locality; (6) Credit-g~\cite{kadra2021well}, for distinguishing good versus bad credit risks; (7) Diabetes\footnote{\texttt{kaggle.com/datasets/uciml/pima-indians-diabetes-database}}, which determines the presence or absence of diabetes in individuals; (8) Heart\footnote{\texttt{kaggle.com/datasets/fedesoriano/heart-failure-prediction}}, which predicts occurrences of coronary artery disease; and (9) Myocardial~\cite{misc_myocardial_infarction_complications_579}, which identifies individuals with chronic heart failure.

We further expand our evaluation with newer and synthetic datasets that are infrequently discussed in literature and were not part of the LLMs’ pretraining data: (10) Cultivars, for forecasting high-yield grain production from soybean varieties\footnote{\texttt{archive.ics.uci.edu/dataset/913}}; (11) NHANES, for classifying the age group from personal records\footnote{\texttt{archive.ics.uci.edu/dataset/887}}; (12) Sequence-type, which categorizes sequences as arithmetic, geometric, Fibonacci, or Collatz types; (13) Solution-mix, aimed at predicting if a mixture composed from four solutions has a concentration percentage greater than 0.5. The Cultivars and NHANES datasets have been recently added to the UCI Machine Learning Repository (after September 2023), guaranteeing exclusion from the training data of the LLM API (gpt-3.5-turbo-0613) employed in our model. The last two datasets are synthetically created, eliminating the chance of memorization by the LLM API and allowing for unbiased evaluation.

The datasets differ in their dimensionality and structural complexity, as provided in Table~\ref{Tab:basic-desc}. Every dataset includes explicitly defined attributes with clear semantics. Particularly, datasets such as `Communities' and `Myocardial' present high-dimensional features with over 100 variables, which challenge the LLM's capacity due to its limited input context window.

\vspace*{-0.12in}
\paragraph{Baselines.}
We benchmark our framework against nine competing methods. The first group comprises three traditional tabular modeling techniques: (1) Logistic regression (LogReg); (2) XGBoost~\cite{chen2016xgboost}; and (3) RandomForest~\cite{ho1995random}. The subsequent two baselines presume access to ample unlabeled data, and include (4) SCARF~\cite{bahri2022scarf}, and (5) STUNT~\cite{nam2023stunt}. It should be acknowledged that, in practice, collecting extensive amounts of unlabeled data can be challenging. Another recent approach, (6) TabPFN~\cite{hollmann2022tabpfn}, leverages synthetic data distributions for model pretraining. The last three baseline methods are LLM-based: (7) In-context learning (In-context)~\cite{wei2022emergent}, (8) TABLET~\cite{slack2023tablet}, and (9) TabLLM~\cite{hegselmann2023tabllm}. The in-context learning approach inserts several labeled examples into the prompt without any further training. TABLET supplements the prompt with additional contextual signals derived from an auxiliary classifier using rule-sets and prototype representations to improve inferential outcomes. TabLLM utilizes parameter-efficient adaptation methods such as IA3~\cite{liu2022few} to fine-tune the LLM using a limited number of labeled samples.

\vspace*{-0.12in}
\paragraph{Implementation details.}
\textsf{FeatLLM} is constructed atop GPT-3.5, though it is inherently flexible with respect to the LLM backbone employed (see Appendix~\ref{sec:backbones} for comparative results with alternative backbones). For LLM inference, we apply a temperature parameter of 0.5 and retain the API default of 1 for top-p. The ensemble size is set to 20, and 10 rules are extracted during feature identification. Additional analysis of hyper-parameter effects is presented in Figure~\ref{fig:hyperparameter2} and Appendix~\ref{sec:hyper-parameter}. The linear model leverages the Adam optimizer with a learning rate fixed at 0.01 and is trained over 200 epochs. We adopt $k$-fold cross-validation to identify the optimal training epoch. When re-implementing baseline approaches, GPT-3.5 is used for in-context learning, whereas TabLLM incorporates T0~\cite{sanh2022multitask} as prescribed in its original protocol. Notably, TabLLM is limited to large language models for which checkpoints are openly released; as a result, GPT-3.5 cannot be utilized with TabLLM. For classical machine learning baselines such as LogReg, XGBoost, and RandomForest, we tune model hyper-parameters using grid search in conjunction with $k$-fold cross-validation. The parameter $k$ is assigned a value of 2 or 4, selected to ensure each training split comprises at least one instance from every class. Further implementation specifics can be found in Appendix~\ref{sec:baselines}.

\vspace*{-0.12in}
\paragraph{Results.}
Tables~\ref{tab:main1} and \ref{tab:main2} provide a detailed comparison of model performances across various datasets. For reliability, each experiment is conducted three times with different random splits of the training and test sets, and both the mean and standard deviation of the AUC scores are reported. \textsf{FeatLLM} either achieves the best results or is ranked second among all baseline methods tested. Remarkably, with only 4-shot samples, our framework matches the performance levels of traditional tabular baselines that necessitate 32-shot samples (refer to Fig.~\ref{fig:compares}). Although STUNT and TabPFN yielded strong improvements on particular datasets, their average performance gains over standard machine learning models were limited. Furthermore, models based on LLM, such as in-context learning, TABLET, and TabLLM, encountered difficulties in processing tabular datasets with many features and could not make inferences. By integrating both bagging and ensemble methodologies, our framework remains robust and effective, even on datasets characterized by a high dimensional feature space. Notably, the concepts of bagging and ensembling could potentially be extended to other LLM-based strategies. However, this would entail a doubling of inference operations per sample, leading to increased computational demands and rendering such adaptations infeasible in practice. 

\subsection{Analysis \& Discussion}
\label{sec:analysis}
\paragraph{Ablation study.} We investigate the individual contributions of key architectural elements through a series of ablation experiments, where we evaluate the effect of: (1) \textsf{-Tuning}: removing the weight tuning procedure in the linear model and instead aggregating binary feature values to determine the logit; (2) \textsf{-Ensemble}: eliminating the ensembling stage; (3) \textsf{-Description}: discarding feature descriptions; and (4) \textsf{-Reasoning}: skipping the Step-1 reasoning instruction when generating rule sets.

Table~\ref{tab:ablation} presents the mean change in AUC across datasets due to each ablation. In every scenario, omitting a component leads to diminished performance. Notably, the improvement attributable to weight tuning increases with more training samples, suggesting that larger datasets enable more accurate rule importance estimation, thus amplifying the impact of tuning. Conversely, incorporating feature descriptions and explicit reasoning instructions is particularly beneficial when few training shots are available, implying that leveraging the LLM's inherent prior knowledge becomes more critical in data-scarce conditions. Finally, ensembling reliably enhances performance across all tested configurations.

\vspace*{-0.12in}
\paragraph{Training \& inference time comparison.} We evaluate and contrast the durations required for both training and inference across several methods. All tests utilize the Adult dataset, with a training set size of 16 shots. Unless otherwise stated, a single A100 GPU is used; specifically, TabLLM relies on four A100 GPUs to support model parallelism. Table~\ref{tab:cost} details the recorded runtimes. Remarkably, \textsf{FeatLLM} demonstrates an inference time that is relatively low and on par with standard baseline approaches. In comparison, other approaches leveraging LLMs—including In-context learning, TABLET, and TabLLM—suffer from substantially increased inference times. With respect to training, \textsf{FeatLLM} performs similarly to the other LLM-driven models, requiring only 30 API calls, which imposes minimal budgetary impact and allows for parallel execution.

\vspace*{-0.12in}
\paragraph{Quality of features generated by \textsf{FeatLLM}.}
We designed a straightforward experiment to assess how informative the rule-based features produced by \textsf{FeatLLM} are for real-world problem-solving. To do this, we calculate the average AUROC between these new features and the respective target labels, followed by determining the fraction of features with AUROC above 0.5 for every dataset. Table~\ref{tab:quality} below presents a summary of these findings. The outcomes reveal that the bulk of the generated rules have strong relevance to the prediction target and yield useful information for the learning task (see the ``Before tuning'' column). Additionally, during the process of fitting the linear model, rules with minimal association to the data (those with weights below a specified threshold) are systematically removed (see the ``After tuning'' column). The threshold was set at 0.1, informed by inspecting the weight histogram.

\vspace*{-0.12in}
\paragraph{Effect of adding spurious correlations.} To evaluate how well different approaches handle the presence of spurious correlations in settings with limited labeled data, we performed a targeted analysis. Specifically, we focused on the Heart dataset, to which we randomly appended features from the Adult dataset that are extraneous to the predictive task of the Heart dataset. We assessed model performance as the number of these irrelevant features was incrementally increased. For consistency, we fixed the number of labeled examples to 8 shots and omitted any use of unlabeled data, thus leaving out methods such as SCARF and STUNT from this particular evaluation.

Figure~\ref{fig:spurious} depicts how performance is influenced by the introduction of spurious correlations. Among the compared methods, \textsf{FeatLLM} exhibited the smallest decline in accuracy as more irrelevant features were added. This resilience can likely be attributed to the way our framework leverages prior knowledge to emphasize the alignment between the primary task and relevant features when generating rules, thereby reducing the risk of forming rules centered on extraneous attributes and enhancing robustness. Supporting this, we find that rules involving noisy features contributed to only 1–2\% of the overall rule set. Nevertheless, it is also evident that \textsf{FeatLLM} is susceptible to assimilating spurious correlations and misconstruing causal relationships when features sourced from data that are highly similar to the target are included (refer to Appendix~\ref{sec:spurious-appendix}).

\vspace*{-0.12in}
\paragraph{Hyper-parameter analysis}
There are two primary hyper-parameters in \textsf{FeatLLM}: the ensemble size and the number of rules extracted per class per LLM query. To assess their effects, we performed empirical analysis varying each. Our findings indicate that increasing the ensemble size enhances performance, but after reaching a threshold (approximately 20 ensembles), improvements plateau (refer to Fig.~\ref{fig:appendix-hyperparameter1} in Appendix). Regarding the number of rules, we observed no statistically significant variations across tested settings, but found that using 10 rules offers an optimal compromise between prompt size and predictive accuracy (see Fig.~\ref{fig:hyperparameter2}). We attribute this to the observation that increasing the rule count expands the dimensionality of the input, which can convey more information but also introduces a greater risk of overfitting, especially in the low-shot context.

\section{Conclusion}
We introduce \textsf{FeatLLM}, an approach that significantly elevates the performance of LLM-driven few-shot learning on tabular datasets. Unlike prior methods focused solely on utilizing LLMs for direct instance-level predictions, \textsf{FeatLLM} leverages LLMs as feature engineers, generating rules to guide prediction and employing an ensemble framework with bagging. This approach incurs minimal inference overhead, eliminates the need for model training by relying solely on API interaction, and bypasses restrictions tied to feature dimensionality. \textsf{FeatLLM} delivers strong prediction capabilities, making it an excellent choice for practical deployment of LLMs on tabular data with limited labeled samples.

While \textsf{FeatLLM} is currently optimized for low-shot learning scenarios, we aim to enhance its applicability to larger datasets, explore advanced feature creation mechanisms beyond rules, and improve interpretability to facilitate broader task applicability in future research.

\section{Broader Impact} 
The \textsf{FeatLLM} framework is designed to harness the pre-existing knowledge embedded in LLMs for superior performance on tabular learning tasks, exceeding what traditional supervised learning approaches offer. By focusing on data-efficient methodologies, our approach makes a significant contribution toward addressing overfitting that arises from small training sets, leveraging external knowledge to compensate. \textsf{FeatLLM} is particularly advantageous in professional environments—such as finance and healthcare—where acquiring labels is often prohibitively expensive or difficult, and where domain-specific information within LLMs provides practical value, given these models are pre-trained on relevant textual data. Looking ahead, examining the influence of potential societal biases or misinformation encoded within LLM prior knowledge is crucial, as such information can heavily influence model decisions, potentially surpassing the influence of observed labeled examples, and will be important for the broader adoption of these systems.

\end{document}